% Options for packages loaded elsewhere
% Options for packages loaded elsewhere
\PassOptionsToPackage{unicode}{hyperref}
\PassOptionsToPackage{hyphens}{url}
\PassOptionsToPackage{dvipsnames,svgnames,x11names}{xcolor}
%
\documentclass[
]{report}
\usepackage{xcolor}
\usepackage[top=30mm,left=20mm,right=20mm,bottom=30mm]{geometry}
\usepackage{amsmath,amssymb}
\setcounter{secnumdepth}{5}
\usepackage{iftex}
\ifPDFTeX
  \usepackage[T1]{fontenc}
  \usepackage[utf8]{inputenc}
  \usepackage{textcomp} % provide euro and other symbols
\else % if luatex or xetex
  \usepackage{unicode-math} % this also loads fontspec
  \defaultfontfeatures{Scale=MatchLowercase}
  \defaultfontfeatures[\rmfamily]{Ligatures=TeX,Scale=1}
\fi
\usepackage{lmodern}
\ifPDFTeX\else
  % xetex/luatex font selection
\fi
% Use upquote if available, for straight quotes in verbatim environments
\IfFileExists{upquote.sty}{\usepackage{upquote}}{}
\IfFileExists{microtype.sty}{% use microtype if available
  \usepackage[]{microtype}
  \UseMicrotypeSet[protrusion]{basicmath} % disable protrusion for tt fonts
}{}
\makeatletter
\@ifundefined{KOMAClassName}{% if non-KOMA class
  \IfFileExists{parskip.sty}{%
    \usepackage{parskip}
  }{% else
    \setlength{\parindent}{0pt}
    \setlength{\parskip}{6pt plus 2pt minus 1pt}}
}{% if KOMA class
  \KOMAoptions{parskip=half}}
\makeatother
% Make \paragraph and \subparagraph free-standing
\makeatletter
\ifx\paragraph\undefined\else
  \let\oldparagraph\paragraph
  \renewcommand{\paragraph}{
    \@ifstar
      \xxxParagraphStar
      \xxxParagraphNoStar
  }
  \newcommand{\xxxParagraphStar}[1]{\oldparagraph*{#1}\mbox{}}
  \newcommand{\xxxParagraphNoStar}[1]{\oldparagraph{#1}\mbox{}}
\fi
\ifx\subparagraph\undefined\else
  \let\oldsubparagraph\subparagraph
  \renewcommand{\subparagraph}{
    \@ifstar
      \xxxSubParagraphStar
      \xxxSubParagraphNoStar
  }
  \newcommand{\xxxSubParagraphStar}[1]{\oldsubparagraph*{#1}\mbox{}}
  \newcommand{\xxxSubParagraphNoStar}[1]{\oldsubparagraph{#1}\mbox{}}
\fi
\makeatother

\usepackage{color}
\usepackage{fancyvrb}
\newcommand{\VerbBar}{|}
\newcommand{\VERB}{\Verb[commandchars=\\\{\}]}
\DefineVerbatimEnvironment{Highlighting}{Verbatim}{commandchars=\\\{\}}
% Add ',fontsize=\small' for more characters per line
\usepackage{framed}
\definecolor{shadecolor}{RGB}{241,243,245}
\newenvironment{Shaded}{\begin{snugshade}}{\end{snugshade}}
\newcommand{\AlertTok}[1]{\textcolor[rgb]{0.68,0.00,0.00}{#1}}
\newcommand{\AnnotationTok}[1]{\textcolor[rgb]{0.37,0.37,0.37}{#1}}
\newcommand{\AttributeTok}[1]{\textcolor[rgb]{0.40,0.45,0.13}{#1}}
\newcommand{\BaseNTok}[1]{\textcolor[rgb]{0.68,0.00,0.00}{#1}}
\newcommand{\BuiltInTok}[1]{\textcolor[rgb]{0.00,0.23,0.31}{#1}}
\newcommand{\CharTok}[1]{\textcolor[rgb]{0.13,0.47,0.30}{#1}}
\newcommand{\CommentTok}[1]{\textcolor[rgb]{0.37,0.37,0.37}{#1}}
\newcommand{\CommentVarTok}[1]{\textcolor[rgb]{0.37,0.37,0.37}{\textit{#1}}}
\newcommand{\ConstantTok}[1]{\textcolor[rgb]{0.56,0.35,0.01}{#1}}
\newcommand{\ControlFlowTok}[1]{\textcolor[rgb]{0.00,0.23,0.31}{\textbf{#1}}}
\newcommand{\DataTypeTok}[1]{\textcolor[rgb]{0.68,0.00,0.00}{#1}}
\newcommand{\DecValTok}[1]{\textcolor[rgb]{0.68,0.00,0.00}{#1}}
\newcommand{\DocumentationTok}[1]{\textcolor[rgb]{0.37,0.37,0.37}{\textit{#1}}}
\newcommand{\ErrorTok}[1]{\textcolor[rgb]{0.68,0.00,0.00}{#1}}
\newcommand{\ExtensionTok}[1]{\textcolor[rgb]{0.00,0.23,0.31}{#1}}
\newcommand{\FloatTok}[1]{\textcolor[rgb]{0.68,0.00,0.00}{#1}}
\newcommand{\FunctionTok}[1]{\textcolor[rgb]{0.28,0.35,0.67}{#1}}
\newcommand{\ImportTok}[1]{\textcolor[rgb]{0.00,0.46,0.62}{#1}}
\newcommand{\InformationTok}[1]{\textcolor[rgb]{0.37,0.37,0.37}{#1}}
\newcommand{\KeywordTok}[1]{\textcolor[rgb]{0.00,0.23,0.31}{\textbf{#1}}}
\newcommand{\NormalTok}[1]{\textcolor[rgb]{0.00,0.23,0.31}{#1}}
\newcommand{\OperatorTok}[1]{\textcolor[rgb]{0.37,0.37,0.37}{#1}}
\newcommand{\OtherTok}[1]{\textcolor[rgb]{0.00,0.23,0.31}{#1}}
\newcommand{\PreprocessorTok}[1]{\textcolor[rgb]{0.68,0.00,0.00}{#1}}
\newcommand{\RegionMarkerTok}[1]{\textcolor[rgb]{0.00,0.23,0.31}{#1}}
\newcommand{\SpecialCharTok}[1]{\textcolor[rgb]{0.37,0.37,0.37}{#1}}
\newcommand{\SpecialStringTok}[1]{\textcolor[rgb]{0.13,0.47,0.30}{#1}}
\newcommand{\StringTok}[1]{\textcolor[rgb]{0.13,0.47,0.30}{#1}}
\newcommand{\VariableTok}[1]{\textcolor[rgb]{0.07,0.07,0.07}{#1}}
\newcommand{\VerbatimStringTok}[1]{\textcolor[rgb]{0.13,0.47,0.30}{#1}}
\newcommand{\WarningTok}[1]{\textcolor[rgb]{0.37,0.37,0.37}{\textit{#1}}}

\usepackage{longtable,booktabs,array}
\usepackage{calc} % for calculating minipage widths
% Correct order of tables after \paragraph or \subparagraph
\usepackage{etoolbox}
\makeatletter
\patchcmd\longtable{\par}{\if@noskipsec\mbox{}\fi\par}{}{}
\makeatother
% Allow footnotes in longtable head/foot
\IfFileExists{footnotehyper.sty}{\usepackage{footnotehyper}}{\usepackage{footnote}}
\makesavenoteenv{longtable}
\usepackage{graphicx}
\makeatletter
\newsavebox\pandoc@box
\newcommand*\pandocbounded[1]{% scales image to fit in text height/width
  \sbox\pandoc@box{#1}%
  \Gscale@div\@tempa{\textheight}{\dimexpr\ht\pandoc@box+\dp\pandoc@box\relax}%
  \Gscale@div\@tempb{\linewidth}{\wd\pandoc@box}%
  \ifdim\@tempb\p@<\@tempa\p@\let\@tempa\@tempb\fi% select the smaller of both
  \ifdim\@tempa\p@<\p@\scalebox{\@tempa}{\usebox\pandoc@box}%
  \else\usebox{\pandoc@box}%
  \fi%
}
% Set default figure placement to htbp
\def\fps@figure{htbp}
\makeatother





\setlength{\emergencystretch}{3em} % prevent overfull lines

\providecommand{\tightlist}{%
  \setlength{\itemsep}{0pt}\setlength{\parskip}{0pt}}



 


\makeatletter
\@ifpackageloaded{caption}{}{\usepackage{caption}}
\AtBeginDocument{%
\ifdefined\contentsname
  \renewcommand*\contentsname{Table of contents}
\else
  \newcommand\contentsname{Table of contents}
\fi
\ifdefined\listfigurename
  \renewcommand*\listfigurename{List of Figures}
\else
  \newcommand\listfigurename{List of Figures}
\fi
\ifdefined\listtablename
  \renewcommand*\listtablename{List of Tables}
\else
  \newcommand\listtablename{List of Tables}
\fi
\ifdefined\figurename
  \renewcommand*\figurename{Figure}
\else
  \newcommand\figurename{Figure}
\fi
\ifdefined\tablename
  \renewcommand*\tablename{Table}
\else
  \newcommand\tablename{Table}
\fi
}
\@ifpackageloaded{float}{}{\usepackage{float}}
\floatstyle{ruled}
\@ifundefined{c@chapter}{\newfloat{codelisting}{h}{lop}}{\newfloat{codelisting}{h}{lop}[chapter]}
\floatname{codelisting}{Listing}
\newcommand*\listoflistings{\listof{codelisting}{List of Listings}}
\makeatother
\makeatletter
\makeatother
\makeatletter
\@ifpackageloaded{caption}{}{\usepackage{caption}}
\@ifpackageloaded{subcaption}{}{\usepackage{subcaption}}
\makeatother
\usepackage{bookmark}
\IfFileExists{xurl.sty}{\usepackage{xurl}}{} % add URL line breaks if available
\urlstyle{same}
\hypersetup{
  colorlinks=true,
  linkcolor={blue},
  filecolor={Maroon},
  citecolor={Blue},
  urlcolor={Blue},
  pdfcreator={LaTeX via pandoc}}


\author{}
\date{}
\begin{document}

\renewcommand*\contentsname{Table of contents}
{
\hypersetup{linkcolor=}
\setcounter{tocdepth}{2}
\tableofcontents
}

\chapter{USER INSTRUCTIONS}\label{user-instructions}

CloudPedagogy Research Object Template\\
v0.1

\begin{center}\rule{0.5\linewidth}{0.5pt}\end{center}

\section{1. Getting Started}\label{getting-started}

\subsection{Install Requirements}\label{install-requirements}

\begin{itemize}
\tightlist
\item
  Install Quarto: https://quarto.org
\item
  (Optional) Install R or Python if you plan to execute code blocks
\end{itemize}

Check installation:

\begin{Shaded}
\begin{Highlighting}[]
\ExtensionTok{quarto}\NormalTok{ check}
\end{Highlighting}
\end{Shaded}

\begin{center}\rule{0.5\linewidth}{0.5pt}\end{center}

\section{2. Explore the Example First}\label{explore-the-example-first}

Before editing the template, render the example project:

\begin{Shaded}
\begin{Highlighting}[]
\BuiltInTok{cd}\NormalTok{ example}
\ExtensionTok{quarto}\NormalTok{ render}
\end{Highlighting}
\end{Shaded}

Open the generated HTML site in \texttt{\_site/}.

This shows a fully populated synthetic healthcare research object,
including governance artefacts and optional AI reflection.

\begin{center}\rule{0.5\linewidth}{0.5pt}\end{center}

\section{3. Creating Your Own Research
Object}\label{creating-your-own-research-object}

\subsection{Step 1 -- Fork or Clone}\label{step-1-fork-or-clone}

Fork this repository or clone it locally.

\subsection{Step 2 -- Work from the Root
Template}\label{step-2-work-from-the-root-template}

The root directory contains the blank template.\\
Edit files under:

\begin{itemize}
\tightlist
\item
  \texttt{/content/}
\item
  \texttt{/analysis/}
\item
  \texttt{/data/}
\item
  \texttt{/governance/}
\end{itemize}

Replace synthetic example text with your own research content.

\subsection{Step 3 -- Define Your Research
Clearly}\label{step-3-define-your-research-clearly}

Ensure you: - State your research question explicitly - Define outcome
and predictors - Describe study design and assumptions - Document
limitations transparently

\begin{center}\rule{0.5\linewidth}{0.5pt}\end{center}

\section{4. Documenting Data Properly}\label{documenting-data-properly}

Complete:

\begin{itemize}
\tightlist
\item
  \texttt{codebook.csv}
\item
  \texttt{data\_dictionary.qmd}
\item
  \texttt{dataset\_card.qmd}
\end{itemize}

You should clearly describe: - Data source - Variable definitions -
Pre-processing steps - Known limitations

If using synthetic data, explicitly state this.

\begin{center}\rule{0.5\linewidth}{0.5pt}\end{center}

\section{5. Completing Governance
Artefacts}\label{completing-governance-artefacts}

The following files are intentionally included:

\begin{itemize}
\tightlist
\item
  \texttt{model\_card.qmd}
\item
  \texttt{risk\_register.md}
\item
  \texttt{decision\_log.md}
\item
  \texttt{reproducibility\_checklist.md}
\end{itemize}

These should be completed before sharing or submitting your work.

Governance artefacts are not optional add-ons; they are part of
responsible research documentation.

\begin{center}\rule{0.5\linewidth}{0.5pt}\end{center}

\section{6. Using the AI Capability Checkpoints
(Optional)}\label{using-the-ai-capability-checkpoints-optional}

The AI Capability Checkpoints are located under \texttt{/ai-support/}.

You should complete them only if AI tools were used in:

\begin{itemize}
\tightlist
\item
  Drafting
\item
  Code generation
\item
  Data synthesis
\item
  Analysis support
\item
  Editing or critique
\end{itemize}

Guidelines:

\begin{itemize}
\tightlist
\item
  Be proportionate.
\item
  Do not overstate AI involvement.
\item
  Emphasise human oversight.
\item
  Maintain accountability with the research team.
\end{itemize}

If no AI tools were used, these sections can remain blank.

\begin{center}\rule{0.5\linewidth}{0.5pt}\end{center}

\section{7. Rendering Your Project}\label{rendering-your-project}

From the project root:

\begin{Shaded}
\begin{Highlighting}[]
\ExtensionTok{quarto}\NormalTok{ render}
\end{Highlighting}
\end{Shaded}

To render a specific folder:

\begin{Shaded}
\begin{Highlighting}[]
\ExtensionTok{quarto}\NormalTok{ render content/04{-}methods.qmd}
\end{Highlighting}
\end{Shaded}

Output will be generated in the \texttt{\_site/} directory.

\begin{center}\rule{0.5\linewidth}{0.5pt}\end{center}

\section{8. Exporting for Sharing}\label{exporting-for-sharing}

Quarto supports:

\begin{itemize}
\tightlist
\item
  HTML output
\item
  PDF output (requires LaTeX/TinyTeX)
\item
  Word export (optional configuration)
\end{itemize}

You may: - Share the HTML site - Export a PDF submission pack - Provide
a repository link for transparency

\begin{center}\rule{0.5\linewidth}{0.5pt}\end{center}

\section{9. Common Pitfalls}\label{common-pitfalls}

\begin{itemize}
\tightlist
\item
  Leaving governance artefacts incomplete\\
\item
  Forgetting to clarify synthetic vs real data\\
\item
  Over-attributing intellectual contribution to AI tools\\
\item
  Mixing template and example content\\
\item
  Failing to render before sharing
\end{itemize}

Always render and review your site before distribution.

\begin{center}\rule{0.5\linewidth}{0.5pt}\end{center}

\section{10. Design Philosophy}\label{design-philosophy}

This template is:

\begin{itemize}
\tightlist
\item
  Lightweight
\item
  Static
\item
  Governance-oriented
\item
  AI-aware (optional)
\item
  Infrastructure-focused
\end{itemize}

It is not a data platform or compliance system.

Its purpose is to structure research transparently and responsibly using
a reproducible framework.

\begin{center}\rule{0.5\linewidth}{0.5pt}\end{center}

\section{Version}\label{version}

v0.1 -- Synthetic Demonstration Release




\end{document}
